\section{The changed command controller and its PA calibration}
\label{sec:pa194c-calibration}

The new machine retains the separate-diagonal 190-token axiom emitter and
changes the controller that uses its output.  It combines two earlier
controller transformations: Shift tests the cached head alone, and
successful modus ponens puts its conclusion in that cache.  The two
transformations must be composed: transplanting only one of their old
proof-word compilers would give incorrect stack identities.

There is a particularly strong source-level cross-check.  After resolving
every structured jump, the complete 3,649-instruction program is identical
to the supplied PA195C program.  The frozen prior instruction array has
SHA-256
\begin{quote}\small
\nolinkurl{0afee1346c0931b3b91fea70571b01fd665a802ab9327409e4030f673412f5e3}.
\end{quote}
The new physical allocation therefore implements an already calibrated
source program.  The proof below also describes that calibration, so that
the change from the existing 3,673-instruction PA194 is explicit.

\subsection{Stores, numerical aliases, and the two changed rules}

Let
\[
 p(a,b)=\frac{(a+b)(a+b+1)}2+b,
 \qquad [a;w]=p(a,w),\qquad I(a,b)=4p(a,b).
\]
Use iterated cons notation $[a_1,\ldots,a_k;w]$, with $a_1$ on top.
The total projections satisfy
$\operatorname{hd}[a;w]=a$, $\operatorname{tl}[a;w]=w$, and
$\operatorname{hd}(0)=\operatorname{tl}(0)=0$.  In particular,
$[0;0]=0$: this encoding does not distinguish an empty stack from a
stack containing only zero entries.  All subsequent identities are
identities of numbers with these total projections.

A command store is $\langle t,w,d,h\rangle$, consisting of top formula,
underlying stack, numerical offset, and cached formula.  The modified
Shift is
\[
 S\langle t,w,d,h\rangle=
 \begin{cases}
  \mathsf{halt},&h=1,\\
  \langle h,\operatorname{tl}(w),d,\operatorname{hd}(w)\rangle,&h\ne1.
 \end{cases}                                                    \tag{C1}
\]
For $y=d+x$, the modified modus-ponens command is
\[
 \mathsf{MP}_x\langle t,w,d,h\rangle=
 \begin{cases}
  \langle0,w,0,y\rangle,&I(t,y)=h,\\
  \langle0,w,y,0\rangle,&I(t,y)\ne h.
 \end{cases}                                                    \tag{C2}
\]
Only Shift can halt.  The original PA194 controller tests $t+h=1$ at
Shift and returns $\langle y,[0;w],y,0\rangle$ on successful MP.
Thus the new controller has different command semantics, rather than
only renamed physical controls.  For example, $I(0,2)=20$, and MP with
operand two sends $\langle0,0,0,20\rangle$ to
$\langle0,0,0,2\rangle$ in the new controller and to
$\langle2,0,2,0\rangle$ in the original one.

Write $U(a,B)=4p(B,a)+2$ and $\Gamma(B)=I(B,U(B,B))$.
For $D=d+x$, $H=h+y$, and $A=c$, the unchanged AX command is
\[
 \mathsf{AX}_{c,x,y}\langle t,w,d,h\rangle
  =\langle K(H,D,c),[q_1,\ldots,q_9,t;w],0,0\rangle,            \tag{C3}
\]
where, from stack top to bottom,
\[
 (q_1,\ldots,q_9)=
 (\Gamma(A),L_6,L_5,L_4,L_3,L_2,L_1,
   \rho(D,A,H,d),U(d,h)).
\]
Here $L_1,\ldots,L_6$ are exactly the six inherited logical heads, and
\[
 \rho(D,A,H,d)=((H\to d)\to D)\to((D\to H)\to(A\to H))
\]
denotes the corresponding numerical implication code, with $d\le D$.
The top $K$ is the inherited separation/vacuity gate, replaced by zero
when its numerical freshness guard fails.  The guard does not suppress
the separately pushed diagonal instance.  Every $q_j$ except $q_9$,
and also $K$, has residue zero modulo four; $q_9$ has residue two.
Consequently none of these generated entries is the false code one.
This is an arithmetic fact about constructors and requires no consistency
assumption.  The complete 190-token emitter word, not just its head set,
is unchanged from PA194.

\subsection{Uniform stack identities for ten pushed entries}

Commands in a word execute from left to right.  Iterating (C1) after (C3)
gives, for every incoming store,
\begin{align*}
 \mathsf{AX}_{c,x,y};S^9:
 \langle t,w,d,h\rangle
  &\longmapsto
   \langle\rho(d+x,c,h+y,d),[t;w],0,U(d,h)\rangle,\\
 \mathsf{AX}_{c,x,y};S^{10}:
 \langle t,w,d,h\rangle
  &\longmapsto \langle U(d,h),w,0,t\rangle.
\end{align*}
The first Shift tests zero; the next nine test generated entries.
None of these ten tests examines the saved top $t$.  These identities
therefore hold even when $t=1$ and at every zero-stack alias.

Define
\begin{align*}
 \mathsf{Push}&=\mathsf{AX}_{0,0,0};S^9;\mathsf{MP}_0,\\
 \mathsf{Pop}&=\mathsf{AX}_{0,0,0};S^{11},\\
 \mathsf{Discard}&=\mathsf{MP}_1;\mathsf{AX}_{1,0,0};S^{10};
                    \mathsf{MP}_0.
\end{align*}
They satisfy
\begin{align*}
 \mathsf{Push}\langle t,w,d,h\rangle
  &=\langle0,[t;w],0,0\rangle,\\
 \mathsf{Pop}\langle t,[a;w],d,h\rangle
  &=\begin{cases}
    \mathsf{halt},&t=1,\\
    \langle t,w,0,a\rangle,&t\ne1,
   \end{cases}\\
 \mathsf{Discard}\langle t,w,0,0\rangle
  &=\langle0,w,0,0\rangle.
\end{align*}
For Push, its last comparison is between an implication code and the
universal code $U(d,h)$, so MP fails with offset zero.  Pop's last Shift
tests the saved top and, if safe, removes the required underlying entry.
For Discard, $I(t,1)>0$ forces its first MP to fail; the AX identity
then gives $\langle U(1,0),w,0,0\rangle$, and
$I(U(1,0),0)>0$ forces its final MP to fail with offset zero.
Pop's possible early halt is intentional and is included in the proof
compiler, rather than excluded by an unproved noncontradiction invariant.

Let $\mathsf{Offset}_0$ be empty and
$\mathsf{Offset}_i=\mathsf{MP}_i$ for $i>0$.  It maps
$\langle0,w,0,0\rangle$ to $\langle0,w,i,0\rangle$.
The word
\[
 \mathsf{Gen}_i=\mathsf{Push};\mathsf{Offset}_i;S;
                         \mathsf{AX}_{i,0,0};S^{10}
\]
maps $\langle a,w,0,0\rangle$ to
$\langle U(i,a),w,0,0\rangle$ for every $a,w$, including $a=1$.
The Shift before AX tests cache zero and fetches $a$ into the cache;
AX clears that cache before it could be tested.

For $1\le j\le7$, the selector
\[
 \mathsf{Discard};\mathsf{AX}_{A,D,H};S^{j+1};\mathsf{Pop}^{9-j}
\]
returns $\langle q_j,w,0,0\rangle$ from
$\langle t,w,0,0\rangle$.  This includes the diagonal head at $j=1$.
Each Pop preserves a generated implication code and hence is safe.
For $u\le a$, the restricted propositional selector is
\[
 \mathsf{Discard};\mathsf{Offset}_u;
       \mathsf{AX}_{b,a-u,c};S^9;\mathsf{Pop},
\]
and the gate selector is
$\mathsf{Discard};\mathsf{AX}_{c,b,a};\mathsf{Pop}^{10}$.
These return respectively $\rho(a,b,c,u)$ and $K(a,b,c)$,
preserving the arbitrary underlying stack and clearing both caches.

The new MP suffix is particularly short:
\[
 \mathsf{MPSuffix}_b=S;S;\mathsf{MP}_b;S.
\]
From $\langle0,[a,I(a,b),0;w],0,0\rangle$, it halts if
$a=1$ or $b=1$, and otherwise returns $\langle b,w,0,0\rangle$.
The first Shift fetches $a$; the second tests it and fetches the major
premise; MP puts $b$ in the cache; the final Shift tests $b$ and consumes
the reserved zero.  Without that zero, the suffix would consume an entry
of the framed stack.  Total Cantor decoding proves the same identity
when $a=b=w=0$.

\subsection{Completeness and soundness at the command level}

Let $\mathcal C$ be the numerical separate-diagonal calculus implemented
by the unchanged emitter, with its inherited restricted propositional
axioms, generalization of theorems, and modus ponens.
For each finite proof tree $p$ of $\mathcal C$, construct a word $W_p$
which, from every $\langle t,w,0,0\rangle$, either halts or returns
$\langle\operatorname{concl}(p),w,0,0\rangle$.
Use the displayed selectors at axiom leaves and
$W_{p_0};\mathsf{Gen}_i$ at a generalization node.  At an MP node with
minor premise $p_0$ proving $a$ and major premise $p_1$ proving $I(a,b)$,
use
\[
 \mathsf{Discard};\mathsf{Push};W_{p_1};\mathsf{Push};
                   W_{p_0};\mathsf{Push};\mathsf{MPSuffix}_b.
\]
The first two words reserve the zero required by the MP suffix.
Induction on the proof tree proves the stated contract.  Every early
halt is allowed, and any suffix appended to an already halting word
preserves that halt.  If $p$ proves one, then
\[
 \mathsf{AX}_{0,0,0};W_p;\mathsf{Pop}
\]
halts from every incoming store, including stores with arbitrary caches:
the prefix AX clears both caches, and the final Pop tests the conclusion
if an earlier command has not already halted.  The inherited scheduler
enumerates every finite command word.  The universal incoming-store
property is essential because previous trials can leave a store behind.

For the converse, annotate the top $t$, cached formula $h$, and every
live canonical stack entry with a $\mathcal C$ proof.  No derivability
condition is imposed on the numerical offset $d$.  For $w>0$ one has
$\operatorname{tl}(w)<w$, so repeated tail projection reaches zero in
at most $w$ steps and gives a primitive-recursive finite canonical stack.
The fixed proof $\rho(0,0,0,0)=0$ supplies the implicit zero tail.
Thus $[0;0]=0$ requires no infinite annotation.

AX saves the annotated top, generalizes the old cached formula, and emits
axioms; its dominance condition is $d\le d+x$.  MP failure returns
the proved code zero, and MP success applies modus ponens to its two
annotated premises.  Shift moves annotated entries, and a halting Shift
exposes an annotated one.  Induction on the finite command trace proves
that every source halt yields a $\mathcal C$ refutation.  The new command
contracts, the selectors, the proof-tree compiler, and all annotation
updates use primitive-recursive operations.  PA verifies their defining
equations and the indicated finite inductions.

\subsection{The inherited logical bridge to PA}

The logical calibration is inherited from the separate-diagonal PA194 and
PA195C analyses\cite{pa195,pa194,pa195c}; its full argument is retained with the reference material.
Its two directions are syntactic proof transformations, not an appeal to
external soundness of PA.  Here are the interfaces used by the present
construction.

For completeness, the diagonal schema
$\Gamma(B)=B\to\forall x_B B$ replaces equality vacuity.  Every variable
label in a formula of code $B$ is smaller than $B$.  To recover reverse
membership rewriting, put
$C=E(a,b)\to(\chi(b)\to\chi(a))$, take the separator property
$\neg C$, and choose its own numerical code as the fresh separator label.
Pointed separation, the retained direct membership rewrites, the equality
witness, and membership-only vacuity imply the reverse rewrite; the
instance $\Gamma(\neg C)$ discharges the separator without presupposing
equality vacuity.  A diagonal instance for $E(a,b)$ gives the neighborhood
equality interface at a fresh binder, and the witness and two-way
membership rewrites transport that binder.  These are the interfaces
needed by the inherited PA-refutation compiler.  All choices of fresh
labels, formula constructions, case splits for coincident labels, and
finite proof concatenations are primitive recursive and PA-verifiable.

For soundness, interpret membership by binary membership of a finite-set
code, equality by arithmetic equality, and quantification by ordinary
first-order quantification.  PA proves finite binary comprehension for
each arithmetic formula by induction on the finite bound.  Substitution
of the formula syntax into that fixed induction-proof template gives a
primitive-recursive proof-producing operation whose correctness PA
verifies; this is not a uniform arithmetic truth predicate.  Pointed
separation is an instance of finite comprehension, and diagonal vacuity
uses the strict code bounds.  The label-collision case $a=b$ is handled
separately, since the inner binder changes both occurrences of that label.
Code one translates to binary self-membership, which PA refutes.  The
constructor alias $I(0,0)=0$ is handled by the fixed proof equivalence of
truth and truth implying truth.

Consequently the inherited transformations give
\[
 \mathrm{PA}\vdash
 \bigl(\exists p\,\operatorname{Prf}_{\mathcal C}(p,1)\bigr)
 \ \longleftrightarrow\
 \bigl(\exists q\,\operatorname{Prf}_{\mathrm{PA}}
                         (q,\ulcorner0=1\urcorner)\bigr).
\]
Combining this with the command-level constructions gives the corresponding
equivalence for halting of the exact 3,649-instruction source program.
The physical verification then supplies every source-instruction simulation
for arbitrary natural register values and arbitrary finite nonempty spare
lists, together with the exact blank-tape startup.  Positive simulation
durations allow finite traces to be composed in both directions.  All
fixed local cases and the generic list-sweep lemmas are formalizable in PA.
Therefore, for the delivered literal table,
\[
 \boxed{\displaystyle
 \mathrm{PA}\vdash
 \neg\operatorname{Halt}(\mathrm{PA194C},\mathrm{blank})
 \ \longleftrightarrow\ \operatorname{Con}(\mathrm{PA}).}
\]
This is an ordinary mathematical proof of a PA-provability statement,
supported by reproducible finite and symbolic certificates.  The package
does not claim a newly elaborated end-to-end Lean or Rocq theorem.

\subsection{What transfers from ZFC and concise iPRA research}

The useful ZFC ideas are the substitution-free numerical calculus,
small register-machine architecture, careful separation of source
correctness from literal tape simulation, and explicit condensed-detachment
certificates for the propositional basis.  The inspected
\texttt{lukasiewicz\_basis.md} records K, S, Peirce, composition and
ex-falso interfaces, while also distinguishing a named-formula completeness
theorem from its numerical and machine bridges.  The same distinction
governs the present proof obligations.

The concise-iPRA work illustrates how a small family of proof rules can
recover larger logical interfaces by derivation, rather than by encoding
each as an independent primitive.  That is relevant to future emitter
redesign, but its theorem is calibration to iPRA.  A compact iPRA proof
checker does not by itself establish the required translations to and
from PA refutations, nor a binary-machine state bound.  The present result
therefore retains the PA-calibrated emitter and obtains a different
194-state construction by combining its allocation with the changed controller.
